\section{Related work and the scope of comparison}\label{sec:section-02}

Machine learning can model nonlinear conditional relationships in asset returns. Gu, Kelly and Xiu compare a broad range of methods for risk-premium prediction and examine the importance of interactions \citep{gu2020}. Their result motivates careful predictive comparisons; it does not supply a performance target for this project. Their universe, prediction horizon, information set and portfolio construction differ from our small daily classification study. We do not copy their reported economic gains into a benchmark table or interpret a different task as a replication.

A simple trend rule is a demanding control for a more elaborate trading architecture because it can capture persistent price behavior without learning a high-dimensional decision boundary. Time-series momentum has been studied across futures markets \citep{momentum2012}. Our rule combines a trailing price-average condition with intermediate-horizon momentum on the existing equity/ETF universe. It is long-only and financed through the same constrained local ledger as the ML policy. It is therefore a project control inspired by trend-following, not a reproduction of the published futures strategy.

Gradient boosting provides a flexible additive model built through successive loss reduction \citep{friedman2001}. We pair a shallow histogram boosting classifier with regularized multinomial logistic regression to contrast nonlinear and linear decision surfaces. The logistic model is both a component and a useful simplicity control. Fixed capacity and chronological fitting make the comparison interpretable. Increasing model count or training duration after seeing poor outer-period results would change the scientific question and enlarge the search history; it would not repair a failed hypothesis.

Probability forecasts require more than class ranking. Proper scoring rules reward distributions that represent the forecaster's uncertainty \citep{gneiting2007}. Log loss penalizes confident wrong answers strongly, whereas accuracy only considers the largest probability. The Brier score measures squared error across the complete class distribution. Temperature scaling is a low-dimensional way to alter sharpness while preserving each row's ordering \citep{guo2017}. Evidence from other application domains does not imply that temperature scaling must improve financial forecasts. We therefore compare raw and calibrated mixtures on identical outer observations and disclose the samples used to fit the temperature.

Dependent observations complicate uncertainty. Adjacent forward-return labels overlap, and symbols respond to common market shocks. Resampling instrument-date rows independently would overstate the effective sample size. Block resampling follows the general motivation of dependent-data bootstrap methods \citep{kunsch1989}, while HAC estimators provide a complementary approach for autocorrelated averages \citep{newey1987}. Our implementation preserves the entire cross section within a sampled date and examines several block lengths. The assumptions remain approximate when the return-generating process shifts over time.

Multiple testing is especially consequential in strategy research. Financial factor evidence illustrates how an extensive search can make conventional single-test thresholds misleading \citep{harvey2016}. Holm adjustment controls a declared family of comparisons \citep{holm1979}, but no adjustment can recover undocumented researcher choices. The Deflated Sharpe Ratio and probability-of-backtest-overfitting literature address aspects of selection risk \citep{bailey2014,bailey2015}. We do not report those estimators here because a complete comparable trial history and their required assumptions are not available. The current registry makes the new comparison family explicit; it does not erase earlier exposure.

Covariance shrinkage is motivated by the instability of estimated risk matrices \citep{ledoit2004}. Our risk rule uses a fixed blend of sample covariance and its diagonal. This is a transparent regularization assumption, not an implementation of the cited paper's optimal shrinkage estimator. Exchange timing follows the official NYSE calendar \citep{nyse2026}. The operational design also recognizes that a scheduled workflow file does not guarantee execution; workflow failures require direct evidence \citep{github2026}. These references establish methods and constraints. All empirical performance claims below come from this project's saved or newly executed experiments.
